% ------------------------------------------------------------------------------
% Plantilla de tesis - Universidad de Antioquia
% Copyright (c) 2026 Simón Zuluaga Henao
%
% Este archivo hace parte de una adaptación de la "Plantilla (LaTeX) IEEE"
% del Sistema de Bibliotecas de la Universidad de Antioquia.
%
% Distribuido bajo la licencia CC BY-NC-SA 4.0.
% Modificado (2026): metadatos reescritos con expl3 (propiedad + listas).
% ------------------------------------------------------------------------------

% ==================================================
% %%%%% ESTE ARCHIVO NO DEBE SER EDITADO %%%%%
% ==================================================



\ExplSyntaxOn

% ==================================================
% ALMACENAMIENTO
% ==================================================

\prop_new:N \g_thesis_prop          % campos simples: clave -> valor
\seq_new:N  \g_thesis_authors_seq   % elementos: {nombre}{apellido}
\seq_new:N  \g_thesis_advisors_seq  % elementos: {tipo}{nombre}{grado}

% ==================================================
% CAMPOS SIMPLES
% \comando{valor} guarda el valor bajo la clave indicada
% ==================================================

\cs_new_protected:Npn \thesis_field:Nn #1#2
  { \NewDocumentCommand #1 {m} { \prop_gput:Nnn \g_thesis_prop {#2} {##1} } }

\thesis_field:Nn \setLanguage    {language}

\thesis_field:Nn \titleinfo      {title}
\thesis_field:Nn \titleshort     {title-short}
\thesis_field:Nn \worktype       {work-type}
\thesis_field:Nn \degree         {degree}

\thesis_field:Nn \faculty        {faculty}
\thesis_field:Nn \program        {program}
\thesis_field:Nn \city           {city}
\thesis_field:Nn \gradyear       {year}

\thesis_field:Nn \rectorName     {rector}
\thesis_field:Nn \deanName       {dean}
\thesis_field:Nn \departmentHead {department-head}

%% Opcionales
\thesis_field:Nn \cohort         {cohort}
\thesis_field:Nn \researchGroup  {research-group}
\thesis_field:Nn \researchCenter {research-center}
\thesis_field:Nn \libraryName    {library}

% Lectura de un campo: \thesis{title}, \thesis{city}, ...
\NewExpandableDocumentCommand \thesis {m}
  { \prop_item:Nn \g_thesis_prop {#1} }

% Condicional de idioma: \IfThesisLanguage{spanish}{si}{no}
\NewDocumentCommand \IfThesisLanguage {m +m +m}
  { \str_if_eq:eeTF { \prop_item:Nn \g_thesis_prop {language} } {#1} {#2} {#3} }

% ==================================================
% AUTORES Y ORIENTADORES (cualquier cantidad)
% ==================================================

\NewDocumentCommand \thesisauthor {mm}
  { \seq_gput_right:Nn \g_thesis_authors_seq {{#1}{#2}} }

\NewDocumentCommand \thesisadvisor {mmm}
  { \seq_gput_right:Nn \g_thesis_advisors_seq {{#1}{#2}{#3}} }

% Formato de cada entrada (modificar aquí si se requiere otro formato)
\cs_new:Npn \__thesis_author:nn   #1#2   { #1~#2 }
\cs_new:Npn \__thesis_advisor:nnn #1#2#3 { #1 \\ #2,~#3 }

% \thesisauthors[separador]   (separador por defecto: \par)
\NewDocumentCommand \thesisauthors { O{\par} }
  {
    \seq_set_map:NNn \l_tmpa_seq \g_thesis_authors_seq { \__thesis_author:nn ##1 }
    \seq_use:Nn \l_tmpa_seq {#1}
  }

% \thesisadvisors[separador]  (separador por defecto: \par\medskip)
\NewDocumentCommand \thesisadvisors { O{\par\medskip} }
  {
    \seq_set_map:NNn \l_tmpa_seq \g_thesis_advisors_seq { \__thesis_advisor:nnn ##1 }
    \seq_use:Nn \l_tmpa_seq {#1}
  }

% ==================================================
% CITA (formato IEEE para la página legal)
% ==================================================

\seq_new:N \l__thesis_names_seq
\seq_new:N \l__thesis_initials_seq
\tl_new:N  \l__thesis_first_tl

% " y " / " and " según el idioma (los espacios van dentro de cada rama)
\cs_new:Npn \__thesis_and:
  { \str_if_eq:eeTF { \prop_item:Nn \g_thesis_prop {language} } {spanish} {~y~} {~and~} }
\cs_new:Npn \__thesis_last_sep:
  { \str_if_eq:eeTF { \prop_item:Nn \g_thesis_prop {language} } {spanish} {~y~} {,~and~} }

% Primera letra de una palabra (funciona con tildes: Á, É, Ñ...)
\cs_new_protected:Npn \__thesis_first_char:n #1
  { \text_map_inline:nn {#1} { ##1 \text_map_break: } }

\cs_new_protected:Npn \__thesis_initials:n #1
  {
    \seq_set_split:Nnn \l__thesis_names_seq { ~ } {#1}
    \seq_set_map:NNn \l__thesis_initials_seq \l__thesis_names_seq
      { \__thesis_first_char:n {##1} . }
    \seq_use:Nn \l__thesis_initials_seq { ~ }
  }

% {nombre}{apellido} -> "Apellido, N."
\cs_new_protected:Npn \__thesis_author_cite:nn #1#2
  { #2,~\__thesis_initials:n {#1} }

% Lista de autores para la referencia.
\NewDocumentCommand \thesisauthorscite {}
  {
    \seq_set_map:NNn \l_tmpa_seq \g_thesis_authors_seq { \__thesis_author_cite:nn ##1 }
    \seq_use:Nnnn \l_tmpa_seq { \__thesis_and: } { ,~ } { \__thesis_last_sep: }
  }

% Apellido(s) para la cita corta.
\NewDocumentCommand \thesisauthorshort {}
  {
    \seq_if_empty:NF \g_thesis_authors_seq
      {
        \seq_get_left:NN \g_thesis_authors_seq \l__thesis_first_tl
        \exp_last_unbraced:NV \use_ii:nn \l__thesis_first_tl
        \int_compare:nNnT { \seq_count:N \g_thesis_authors_seq } > { 1 }
            { ~et~al. }
      }
  }

\ExplSyntaxOff

% ==================================================
% ADICIONALES
% ==================================================

\newcommand{\palabrasClave}[1]{{\null\noindent\bfseries{\itshape Palabras clave} --- #1}}
\newcommand{\keywords}[1]{{\null\noindent\bfseries{\itshape Keywords} --- #1}}
